\subsection{Run 82, concluded: the moment programme and the closure of the
scale identification}\label{sub:r82e}

The closed form of Theorem \ref{thm:jacobi} turns every moment of the family
into a combinatorial identity.  Write $e_k$ for the elementary symmetric
functions and $X=\FF_N\setminus\{0,1\}$.

\begin{theorem}[family second moments]\label{thm:momentpattern}
For every $k\ge2$ and $q\neq0$,
\[
  M_{2k}(q):=\frac{1}{(N-1)^2}\sum_{\chi,\psi\in\widehat G}|J(\chi,\psi;q)|^{2k}
  \;=\;\sum_{(P,Q)\in\FF_N^2}|W_k(P,Q)|^2,
\]
where
\[
  W_k(P,Q)=\sum_{\substack{y\in X^k\\ e_k(y)=P,\; e_k(y-1)=Q}}\ep\!\bigl(q\textstyle\sum_i y_i\bigr)
\]
is the trace of the linear phase $\ep(q\,e_1)$ over the variety
$V_{P,Q}=\{y:e_k(y)=P,\ e_k(y-1)=Q\}$ of dimension $k-2$; hence
$|W_k(P,Q)|=O(N^{(k-2)/2})$ by Deligne, over $\sim N^2$ fibres.  The
normalised second moments converge to the Catalan numbers:
\[
  \frac{1}{N^k}\sum_{(P,Q)}|W_k(P,Q)|^2\;\longrightarrow\;C_k,
  \qquad C_2=2,\; C_3=5,\; C_4=14.
\]
\end{theorem}

The identity is verified exactly (to $\le1.1\cdot10^{-15}$) at $k=2,3,4$ and
ten primes by two independent algorithms, and a third route (the
$(P,Q,e_1)$-parametrisation) confirms $M_6$ to $2.7\cdot10^{-16}$.

\emph{The constant at $k=3$ is elementary.}  Expanding the indicators of the
two product constraints by additive characters reduces the family second
moment to
\[
  \sum_{(P,Q)}|W_3(P,Q)|^2=\frac1{N^2}\sum_{\alpha,\beta}
  |S(\alpha,\beta)|^2,\qquad
  S(\alpha,\beta)=\sum_{u,v} \ep\bigl(q(u+v+2)-\beta uv\bigr)\,f(c),
\]
with $c=(\alpha+\beta)uv+\alpha(u+v)+\alpha+q$ and $f(c)=(N-2)$ on $c=0$,
$-1-\ep(c)$ elsewhere.  The mass concentrates ($\sim99\%$) on
$\alpha,\beta,\alpha+\beta\neq0$, where the $f=(N-2)$ part lives on the conic
$c=0$ --- a hyperbola --- and its phase sum is a Kloosterman sum
$\Kl(A,B)$: $|S|^2\sim(N-2)^2|\Kl|^2$.  Decomposing by the degeneracy
condition $(\alpha+q)(\beta+q)=q^2$ (equivalently $A=0$ or $\kappa=0$):
\begin{itemize}
\item generic pairs: $\operatorname{avg}|\Kl|^2=N$ by the classical
Kloosterman orthogonality (the two-to-one conic map cancelling the half of
the square condition) --- contributes $1$ in units of $N^5$;
\item the degenerate locus ($\approx N$ pairs): $A=\kappa=0$, the conic
splits into two lines with $2N-1$ points, $|S|^2\sim(N-2)^2(2N-1)^2\sim4N^4$
per point --- contributes $4$.
\end{itemize}
Hence $C_3=5=1+4$ by elementary means: exact identities, the classical
Kloosterman bound and orthogonality, and the count of the degenerate locus.
Measured: generic $1.00/0.98/1.01$, degenerate $3.10/3.36/3.53$, total
$4.10/4.34/4.53\to5$ at $N=37/53/73$.

\emph{The boundary of the elementary route.}  At $k=4$ the dominant part of
$S_4$ is the trace over the cubic surface $c=0$, and its family second moment
is $\operatorname{avg}|T|^2/N^2=12.40\to14$ over all $2652$ parameter pairs at
$N=53$ --- with a \emph{smooth} distribution: $90.3\%$ of the mass in one
bucket, zero points above $50$ per unit.  No degenerate locus analogous to
the $k=3$ split conic exists: the constant $C_4=14$ is the bulk second moment
of a family of cubic surfaces, a monodromy quantity in the sense of Katz.
The hierarchy of mechanisms is thus:
$k=2$ elementary (the two orderings); $k=3$ elementary ($1+4$); $k\ge4$ and
the local law require the same single input --- quantitative equidistribution
(Katz-type second moments) for the surface/entry families.

\emph{The scale identification is parameter-free.}  Three pieces close the
conjugation completely.  (i) The half-index trace
$s=a^{(K-1)/2}+a^{-(K-1)/2}=L_{(K-1)/2}\bmod N$ satisfies
$s^2=a^{K-1}+2+a^{1-K}=t+2=5=t^2-4$ --- one line --- so the branch scale is
$c=\pm L_{(K-1)/2}$ with $c^2=t^2-4$ (verified at six branch values across
three designs and end-to-end at five).  (ii) The metaplectic twist
$\zeta$ of the kernel quantisation, defined by
$\lambda_0=\zeta\,\chi_{k_0}(a)^{-1}$ with $\lambda_0$ the principal $K$-th
root of the torsion scalar $\chi$, is the \emph{unique} fourth root of unity
satisfying $\zeta^K=\chi$ (five designs: $\zeta=i,-i,-1,-1,-1$).  (iii) The
block index $k_0$ is the unique class modulo $K$ solving
$\chi_{k_0}(a)=\zeta/\lambda_0$ (five designs: $\{1\},\{1\},\{8\},\{9\},
\{19\}$).  The two measured ``branches'' $(k_0,+c)$, $(-k_0,-c)$ are one
branch related by the exact model symmetry
$\operatorname{model}(k_0,q')=\operatorname{model}(-k_0,-q')$ (verified to
$10^{-9}$ over all probes).  The counting half of Theorem \ref{thm:torus} is
machine-checked in Lean ($\texttt{formal/Bgs/ScaleLaw.lean}$: the scale-law
instantiation at $p=199$ closed by the kernel with no axioms).

\emph{The inert branch.}  Theorem \ref{thm:torus} extends verbatim to
inert $N$ (verified at $N=1103$, $K=48$): every nonzero level set of $Q$ has
$N+1$ points, the cone is $\{0\}$ (anisotropic norm form), one level set
carries $(N+1)/K=m$ orbits, full spectra agree across six distinct orbits in
each of six level sets to $6.6\cdot10^{-14}$, and the explicit torus
elements exist ($10/10$).  The inert analogue of Theorem
\ref{thm:jacobi} --- the matrix elements over the non-split torus, whose
Hecke basis comes from the deleted representation --- remains open.
